\newpage
\section*{\textsc{General Introduction}}
\addcontentsline{toc}{section}{General Introduction}
\pagenumbering{arabic}
\rhead{\itshape General Introduction}
\lhead{}
\cfoot{\thepage}
\pagestyle{fancy}
\renewcommand{\headrulewidth}{0.4pt}

\begin{flushright}
\textit{``The art of teaching is the art of assisting discovery.''\\
Mark van Doren}
\end{flushright}

\subsection*{Background and Motivations}

The construction of university examinations is a cognitively demanding task.
Instructors must simultaneously satisfy pedagogical objectives, calibrate difficulty
across Bloom's cognitive levels \cite{bloom2001}, maintain linguistic consistency,
and avoid repeating questions from prior academic years.
In Algerian universities, this burden is compounded by a bilingual teaching environment
in which lecture materials are distributed in both French and Arabic, and historical
examination archives exist exclusively as unstructured PDF documents with no
machine-readable indexing.

This project develops \textbf{ExamGen}, an AI-powered exam generation platform
that assists teachers by automatically producing high-quality, pedagogically
calibrated exam questions from their own course materials and historical archives.
The AI acts as a \textit{decision-support tool}: instructors review, edit, and
finalise all generated content before use.

Recent advances in Large Language Models (LLMs) have made automated question
generation (AQG) practical. \citet{isley2025quality} demonstrate (91 courses,
$\approx$1,700 students) that AI-generated items match expert quality when grounded
in course materials. \citet{mucciaccia2025mcq} show that multi-stage prompt
engineering with automated LLM review produces high-quality MCQs without fine-tuning.
\citet{morris2024aqg} introduce \textit{answer extractability} as a core quality
metric. \citet{meissner2024itemforge} confirm the need for independent cognitive-level
verification of LLM-generated items.

\subsection*{Problematic and Contributions}

Despite these advances, no existing system simultaneously addresses:
\begin{enumerate}
  \item Bilingual document corpora (French and Arabic) from heterogeneous PDF sources;
  \item Intelligent Document Processing (IDP) including VLM transcription of embedded diagrams;
  \item Language-aware hybrid RAG routing queries to Arabic-specific or multilingual embeddings;
  \item Schema-enforced generation with an independent post-generation quality evaluator.
\end{enumerate}

This report presents \textbf{ExamGen}, an end-to-end system that addresses all four
requirements. The system was developed iteratively: an initial prototype revealed
three concrete failures (absent structured-output enforcement, low free-tier
instruction-following fidelity, and topic-blending from semantic retrieval), each
directly addressed by findings from the reference literature.

\subsection*{Thesis Statement}

This report argues that a bilingual, schema-enforced, RAG-augmented pipeline ---
combining intelligent document processing, language-aware hybrid retrieval, and an
independent post-generation quality evaluator --- can produce examination questions
that are factually grounded, cognitively calibrated, and stylistically consistent
with an instructor's historical practice, without requiring model fine-tuning.
ExamGen demonstrates this thesis on a real corpus of 137 bilingual (French--Arabic)
university examination documents, achieving 87.5\% factual accuracy, 83.8\%
retrieval effectiveness, and reducing instructor preparation time through a
multi-tenant web platform with real-time SSE streaming.

\subsection*{Report Organisation}

Chapter~\ref{chap:soa} reviews related work.
Chapter~\ref{chap:arch} documents system design and architecture.
Chapter~\ref{chap:impl} details the implementation and key algorithms.
Chapter~\ref{chap:results} presents quantitative results.
A general conclusion closes the report.
